\section*{第 \ref{ch:g4:how-sound-travels} 章——声音是如何传播的}

\begin{solution}{exo:g4:how-sound-travels:1}
传播介质：\omterm{def:g2:air-around-us:air}{空气}、水、\omterm{def:g3:solids-liquids-gases:solid}{固体}（木头、砖块、钢铁、绷紧的绳子）。
完全没有介质：真空——像没有\omterm{def:g2:air-around-us:air}{空气}的\omterm{def:g2:moon-in-the-sky:moon}{月球}，或
抽成真空的玻璃罐。
\end{solution}

\begin{solution}{exo:g4:how-sound-travels:2}
\omterm{def:g2:air-around-us:air}{空气}——\omterm{def:g3:sound-around-us:sound}{声音}的介质——被抽走了。
\omterm{def:g3:sound-around-us:vibration}{振动}的铃铛已没有可以带动的介质，
这说明\omterm{def:g3:sound-around-us:sound}{声音}只能通过物质传播，不能穿过虚无。
\end{solution}

\begin{solution}{exo:g4:how-sound-travels:3}
没有介质就没有\omterm{def:g3:sound-around-us:sound}{声音}：\omterm{def:g2:moon-in-the-sky:moon}{月球}上没有\omterm{def:g2:air-around-us:air}{空气}，
大声喊叫时产生的\omterm{def:g3:sound-around-us:vibration}{振动}无从传播，在嘴边就消失了。
\end{solution}

\begin{solution}{exo:g4:how-sound-travels:4}
绷紧的细绳负责传递\omterm{def:g3:sound-around-us:sound}{声音}：说话使杯子\omterm{def:g3:sound-around-us:vibration}{振动}，
杯子带动细绳\omterm{def:g3:sound-around-us:vibration}{振动}，细绳再把\omterm{def:g3:sound-around-us:sound}{声音}送到远处的杯子里。
要点：细绳要拉紧，
且不能碰到其他东西（不能松弛，也不能被手指捏住）。
\end{solution}

\begin{solution}{exo:g4:how-sound-travels:5}
\omterm{def:g3:sound-around-us:sound}{声音}每秒约传播 \qty{340}{m}，大约三秒传播一千米。
\end{solution}

\begin{solution}{exo:g4:how-sound-travels:6}
$9 \div 3 = 3$：约 \qty{3}{km} 远。
三秒后听到雷声，意味着距离约为 \qty{1}{km}：暴风雨已经近多了，快进屋。
\end{solution}

\begin{solution}{exo:g4:how-sound-travels:7}
木头比\omterm{def:g2:air-around-us:air}{空气}更能传声：轻微的刮擦声，
在\omterm{def:g2:air-around-us:air}{空气}中难以传过桌面，
但\omterm{def:g3:solids-liquids-gases:solid}{固体}的桌面能把\omterm{def:g3:sound-around-us:vibration}{振动}直接清晰地送进你的耳朵。
\end{solution}

\begin{solution}{exo:g4:how-sound-travels:8}
\omterm{ex:g4:how-sound-travels:echo}{回声}是你自己的\omterm{def:g3:sound-around-us:sound}{声音}碰到大而硬的墙壁后反射回来，
且返回得足够晚，能被单独听到。
悬崖越远，\omterm{def:g3:sound-around-us:sound}{声音}往返路程越长，\omterm{ex:g4:how-sound-travels:echo}{回声}就来得越晚。
\end{solution}

\begin{solution}{exo:g4:how-sound-travels:9}
船向下发出短促的\omterm{def:g3:sound-around-us:sound}{声音}，并计算来自海底的\omterm{ex:g4:how-sound-travels:echo}{回声}时间：
\omterm{ex:g4:how-sound-travels:echo}{回声}很快返回，说明水浅；\omterm{ex:g4:how-sound-travels:echo}{回声}较晚，说明水深。
海豚在暗水中这样做，蝙蝠在漆黑的山洞里也这样做。
\end{solution}

\begin{solution}{exo:g4:how-sound-travels:10}
第一，钢铁传声远比\omterm{def:g2:air-around-us:air}{空气}好，
所以隆隆声能经长途传播而不消散。
第二，铁轨把\omterm{def:g3:sound-around-us:sound}{声音}沿一条路直接引到侦察兵的耳中，
而在空地上，火车的\omterm{def:g3:sound-around-us:sound}{声音}会像池塘里的波纹向四面扩散，
每走一米都会变弱，
因此铁轨传来的消息远比\omterm{def:g2:air-around-us:air}{空气}中的更响亮、更清楚，早早就到达了。
\end{solution}

\begin{solution}{exo:g4:how-sound-travels:11}
\omterm{def:g3:sound-around-us:sound}{声音}一秒约传播 \qty{340}{m}，
但这一秒包括去程和\emph{返程}。
墙的距离是一半，约为 \qty{170}{m}。
\end{solution}
